\documentclass[10pt]{amsart}
\usepackage{calc,amssymb,amsmath,ulem,stmaryrd,fullpage}
\usepackage{enumerate}
\usepackage{alltt}
\usepackage{multicol}
\RequirePackage[dvipsnames,usenames]{color}
\let\mffont=\sf
\normalem
%\input{kmacros3.sty}
\input{xy}
\xyoption{all}
\usepackage{hyperref}
\usepackage{graphicx}
\usepackage[all,cmtip]{xy}
\usepackage{verbatim}
\usepackage[left=0.95in,top=0.95in,right=0.95in,bottom=0.95in]{geometry}
\newcommand{\tauCohomology}{T}
\newcommand{\FCohomology}{S}


\theoremstyle{theorem}
%\newtheorem*{mainthm*}{Main Theorem}

\newcommand{\note}[1]{\marginpar{\sffamily\tiny #1}}

\def\todo#1{\textcolor{Mahogany}%
{\sffamily\footnotesize\newline{\color{Mahogany}\fbox{\parbox{\textwidth-15pt}{#1}}}\newline}}

\renewcommand{\O}{\mathcal O}
\newcommand{\ie}{{\itshape i.e.} }
\newcommand{\roundup}[1]{\lceil #1 \rceil}
\newcommand{\rounddown}[1]{\lfloor #1 \rfloor}
\renewcommand{\to}[1][]{\xrightarrow{\ #1\ }}
\newcommand{\onto}[1][]{\protect{\xrightarrow{\ #1\ }\hspace{-0.8em}\rightarrow}}
\newcommand{\into}[1][]{\lhook \joinrel \xrightarrow{\ #1\ }}
%\newcommand{DBDual}[1]{{\underline{\omega}_{#1}}}
\newcommand{\DBDual}[1]{{{\uuline \omega}{}^{\mydot}_{#1}}}
\newcommand{\Tt}{{\mathfrak{T}}}
%\newcommand{\bn}{{\mathfrak{n}} }
\newcommand{\sol}[1]{ \vskip 6pt{\color{blue} {\bf Solution:} #1} \vskip 6pt}

\newcommand{\bR}{{\mathbb{R}}}
\newcommand{\va}{{\vec{a}}}
\newcommand{\vb}{{\vec{b}}}
\newcommand{\vzero}{{\vec{0}}}
\newcommand{\vi}{{\vec{i}}}
\newcommand{\vj}{{\vec{j}}}
\newcommand{\vk}{{\vec{k}}}
\newcommand{\vh}{{\vec{h}}}
\newcommand{\vr}{{\vec{r}}}
\newcommand{\vu}{{\vec{u}}}
\newcommand{\vv}{{\vec{v}}}
\newcommand{\vw}{{\vec{w}}}
\newcommand{\vx}{{\vec{x}}}
\newcommand{\vy}{{\vec{y}}}
\newcommand{\vz}{{\vec{z}}}
\newcommand{\vT}{{\vec{T}}}
\newcommand{\vN}{{\vec{N}}}
\newcommand{\vF}{{\vec{F}}}
\newcommand{\vG}{{\vec{G}}}
\newcommand{\bF}{\mathbb{F}}
\newcommand{\bQ}{\mathbb{Q}}
\newcommand{\bZ}{\mathbb{Z}}
\newcommand{\bC}{\mathbb{C}}


\begin{document}
\title{Worksheet \#4-- Math 5320\\Fall 2019}
\author{Due Monday, March 2nd, in Gradescope}

\maketitle
You may work in groups of up to 3 on this assignment.  Only one assignment needs to be turned in per group.  It still needs to be turned in on gradescope.  
\vskip 3pt
First we recall:
\vskip 3pt
\noindent
{\bf Eisenstein's Criterion.}  {\it Suppose $f(x) = a_n x^n + \dots + a_0 \in \bZ[x]$ is a nonconstant polynomial and $p \in \bZ$ is prime.  If $p$ does not divide $a_n$, $p$ divides every other $a_i$ and $p^2$ does not divide $a_0$, then $f(x)$ is irreducible in $\bQ[x]$.}
\vskip 3pt
\noindent
{\bf 1.}  First suppose that $\phi : \bZ[x] \to \bZ[x]$ is a ring homomorphism.  Noticed it is completely determined by where it sends $1$ and suppose that $\phi(x) = h(x)$ with $\deg h(x) > 0$.  Show that if $\phi(f(x)) = f(h(x))$ is irreducible, then $f(x)$ is also irreducible. 
\newpage
\noindent
{\bf 2.}  With notation as in {\bf 1.}, give at least four examples of non-identity ring homomorphisms $\phi$ such that $\phi$ is an isomorphism.
\vskip 150pt
\noindent
{\bf 3.}  Suppose $p \in \bZ$ is prime.  Notice that the polynomial 
\[
    C_p(x) = x^{p-1} + x^{p-2} + \dots + x^{1} + 1 = {x^p - 1 \over x-1} 
\]
is not something that you can apply Eisenstein's criterion to  (this is called the $p$th \emph{cyclotomic polynomial}).  However, in the special cases of $p = 2, 3, 5$, explicitly show that $C(x+1)$ is irreducible by using Eisenstein's criterion.  
\newpage
\noindent
{\bf 4.}  Show that for any prime $p$ that 
\[
    C_p(x+1) = {(x+1)^p - 1 \over x} 
\]
is irreducible.  You may use the binomial formula.  In particular show that ${p \choose i}$ is divisible by $p$ for each $1 \leq i \leq p-1$.
\newpage
\noindent
{\bf 5.}  Show that each of the following polynomials are irreducible.
\vskip 3pt
{\bf (a)}  $x^5 - 4x^2 + 66$.
\vskip 3pt
{\bf (b)}  $x^3 + 3x^2 + 5x + 5$.
\vskip 3pt
{\bf (c)}  $30 x^n - 91$ for any $n > 0$.  
\vskip 3pt
{\bf (d)}  $x^4 + 4x + 1$.
\newpage
\noindent
{\bf 6.}  Consider the polynomial $f(x) = x^3 + x^2 + x + 2$.  In which of the following rings of polynomials is $f$ irreducible?  Justify your answer.
\begin{itemize}
\item[(a)]  $\bR[x]$
\item[(b)]  $\bC[x]$
\item[(c)]  $\bZ/{2}[x]$
\item[(d)]  $\bZ/{3}[x]$
\item[(e)]  $\bZ/{5}[x]$
\item[(f)]  $\bQ[x]$
\end{itemize}
\newpage
\noindent
{\bf 7.}  Consider $3x^2 + 4x + 3 \in \mathbb{Z}_{\text{mod} 5}[x]$.  Show it factors both as $(3x + 2)(x+4)$ and as $(4x + 1)(2x+3)$.  Explain why this \emph{does NOT} contradict unique factorization of polynomials.
\vskip 250pt
\noindent
{\bf 8.}  Show that $x^4 + 1$ is irreducible in $\bQ[x]$ but not irreducible in $\bR[x]$.
\vskip 3pt
\emph{Hint: } For $\bQ[x]$, try a linear shift.  For $\bR[x]$, try a factorization into two quadratic terms


\end{document} 